\subsection{Selection-Backdoor Criterion}
\totalscore{0}
%\score{1}

\begin{figure}[ht]
    \centering
    \begin{tikzpicture}[scale=1, transform shape]
        \node[ndout] (X) at (0,0) {$X$};
        \node[ndout] (Y) at (3,0) {$Y$};
        \node[ndout] (Z) at (1.5,1.5) {$Z$};
        \node[ndout] (S) at (1.5,-1.5) {$S$}; 
        \draw[arout] (X) to (Y);
        \draw[arout] (X) to (S);
        \draw[arout] (Z) to (X);
        \draw[arout] (Z) to (Y);
        \draw[arout] (Z) to (S);
    \end{tikzpicture}
    \caption{Acyclic Directed Mixed Graph (ADMG) $\hat{G}$.}
    \label{fig:cadmg-3}
\end{figure}

Recall that we can recover the causal effect $P(Y|\doit(X))$ using the backdoor criterion for simple SCMs via the adjustment formula $P(Y|\doit(X)) = P(Y|X, Z)\circ P(Z)$ for some set $Z\subseteq V$, if we have the following $\sigma$-separations:
\begin{equation*}
    Z\perp^\sigma I_X \quad \text{and} \quad Y\perp^\sigma I_X|X\cup Z
\end{equation*}

\begin{enumerate}
    \item Draw the graph with the intervention node $I_X$, and apply the backdoor criterion to find an expression for $P(Y|\doit(X))$. \textcolor{red}{PB: perhaps use a more challenging graph, or use different variable names.}
    \answer{%
        If we draw the graph with the intervention node $I_X$:
        \begin{center}
            \begin{tikzpicture}[scale=1, transform shape]
                \node[ndout] (IX) at (0,1.5) {$I_X$};
                \node[ndout] (X) at (0,0) {$X$};
                \node[ndout] (Y) at (3,0) {$Y$};
                \node[ndout] (Z) at (1.5,1.5) {$Z$};
                \node[ndout] (S) at (1.5,-1.5) {$S$}; 
                \draw[arout] (IX) to (X);
                \draw[arout] (X) to (Y);
                \draw[arout] (X) to (S);
                \draw[arout] (Z) to (X);
                \draw[arout] (Z) to (Y);
                \draw[arout] (Z) to (S);
            \end{tikzpicture}
        \end{center}
        then we indeed have the $\sigma$-separations $Z\perp I_X$ and $Y\perp I_X|X, Z$. So we may directly apply the adjustment formula to obtain $P(Y|\doit(X)) = P(Y|X, Z)\circ P(Z)$.
    }
    \points{points?}
    \item Prove the backdoor adjustment formula using rules 2 and 3 of do-calculus, and the law of total probability.
    \answer{
        \begin{align*}
            P(Y|\doit(X)) & = P(Y|Z, \doit(X)) \circ P(Z|\doit(X)) &&\text{law of total probability} \\
            & = P(Y|Z, \doit(X)) \circ P(G) &&\text{Rule 3: $Z\perp I_X$}\\
            & = P(Y|Z, X) \circ P(Z) &&\text{Rule 2: $Y\perp I_X| Z, X$}\\
        \end{align*}
    }
    \points{points?}
\end{enumerate}
Suppose that $S$ is a selection node, i.e., it indicates whether we will see the label $Y$ or not. More precisely, we have access to the distributions $P(X,Z)$ and $P(X, Z, Y|S=1)$.
\begin{enumerate}[resume*]
    \item Slightly adjust the backdoor adjustment formula to find an expression for $P(Y|\doit(X))$ in terms of the distributions that are available to us.
    \answer{%
        Notice that by the Markov property we have the conditional independence $Y\Indep S|X, Z$, so we can include the condition $S=1$ in one of the probabilities:
        $$P(Y|\doit(X)) = P(Y|X, Z, S=1)\circ P(Z)$$
    }
    \points{points?}
\end{enumerate}
